\chapter{Prealbumin (Transthyretin)}

\section*{What Is This Test?}
Prealbumin (transthyretin, TTR) is a small transport protein (55 kDa). The liver makes it. It carries thyroxine and retinol (vitamin A). It is a sensitive marker of visceral protein synthesis \cite{bernstein1996prealbumin}. Its half-life is about 2 days. Albumin lasts much longer (20 days). So prealbumin reflects acute changes in nutritional status and hepatic protein synthesis. Clinicians use it widely to detect malnutrition. They also use it to monitor nutritional repletion in hospitalized patients. It helps assess liver synthetic function \cite{beck2002prealbumin}. Prealbumin can also mark inflammation as a negative acute-phase reactant. It is the precursor protein in hereditary and wild-type transthyretin amyloidosis.

\section*{Alternative Names}
\begin{itemize}
  \item Prealbumin
  \item Transthyretin
  \item TTR
  \item Thyroxine-Binding Prealbumin
  \item TBPA
\end{itemize}
\section*{Principle}
The test uses an immunonephelometric or immunoturbidimetric assay. The antibody is specific for prealbumin. Antigen-antibody complexes scatter light. Scattering is proportional to the prealbumin concentration. A serum separator or EDTA tube is acceptable. Fasting is not required.

\section*{History}
Prealbumin was named for its electrophoretic mobility just ahead of albumin. In the 1960s, it was identified as a carrier protein for thyroxine and retinol. In the 1980s, the hepatology and nutrition communities adopted it. They used it as a short-half-life nutritional marker. In the 1990s, the hereditary amyloidosis community recognized transthyretin mutations. These mutations cause ATTR amyloidosis \cite{liz2020narrative}.

\section*{Why This Test Is Ordered}
\begin{itemize}
  \item Nutritional assessment in hospitalized patients (especially before GI surgery, dialysis, ICU stay)
  \item Monitoring response to enteral or parenteral nutrition
  \item Visceral protein synthesis assessment in liver disease
  \item Identification of protein-energy malnutrition in chronic illnesses, dialysis, elderly
  \item Investigation of older adults with weight loss and frailty
  \item Cardiology workup for patients with restrictive cardiomyopathy (assessment of transthyretin amyloidosis) \cite{conceicao2016redflag}
\end{itemize}

\section*{Primary vs Secondary Test}
Secondary test for nutritional and liver function assessment. It is complementary to albumin.

\section*{How This Test Is Performed}
Venous blood is drawn into a serum separator tube. Results are available within 4--24 hours. Serial measurements every 3--7 days help monitor nutritional repletion.

\section*{What Preparation Is Needed}
None. Fasting is not required.

\section*{Contraindications and Precautions}
\begin{itemize}
  \item Prealbumin is a negative acute-phase reactant. Inflammation reduces values as an unrelated consequence of cytokine-mediated hepatic reprioritization. Concomitant C-reactive protein (CRP) measurement helps interpret. Renal failure increases levels due to reduced clearance. Pregnancy and estrogen therapy raise values.
\end{itemize}
\section*{Risks}
None.

\section*{What Happens After the Test}
Values $<$ 15 mg/dL suggest malnutrition. Values $<$ 10 mg/dL indicate severe malnutrition. An increasing prealbumin over 7--14 days of nutrition support confirms adequate repletion. Interpretation is flagged for inflammation (CRP elevated).

\section*{Understanding Results}

\subsection*{Below Normal}
\begin{itemize}
  \item Mild malnutrition: 10--15 mg/dL
  \item Moderate malnutrition: 5--10 mg/dL
  \item Severe malnutrition: $<$ 5 mg/dL
  \item Liver disease: decreased synthesis (cirrhosis, hepatitis, liver failure)
  \item Acute inflammation: protein redirects to acute-phase reactants (ferritin, CRP)
  \item Hyperthyroidism: increased turnover
  \item Nephrotic syndrome or protein-losing enteropathy
\end{itemize}

\subsection*{Above Normal}
\begin{itemize}
  \item Renal failure (decreased clearance): chronic kidney disease
  \item High-dose corticosteroid therapy
  \item Pregnancy and estrogen therapy: increased hepatic synthesis
  \item Hodgkin lymphoma or HPT hormone use
\end{itemize}

\section*{Normal Results}
\begin{itemize}
  \item \textbf{Adult (serum):} 15--35 mg/dL (150--350 mg/L)
  \item \textbf{Elderly:} slightly lower (15--25 mg/dL)
  \item \textbf{Children:} 12--30 mg/dL
  \item \textbf{Pregnancy (3rd trimester):} typically elevated
  \item \textbf{Nutritional repletion goal:} Rise of $\geq$ 2 mg/dL within 7--14 days suggests adequate nutrition support
\end{itemize}

\section*{Follow-up Tests}
\begin{itemize}
  \item \textbf{Albumin and total protein:} Long-term protein synthesis markers
  \item \textbf{C-reactive protein (CRP):} Confounds acute-phase suppression
  \item \textbf{Comprehensive nutritional assessment:} Anthropometry, dietary intake, micronutrient levels
  \item \textbf{Liver function tests (ALT, AST, INR):} If hepatic synthetic failure suspected
  \item \textbf{Renal function (creatinine, BUN, urinalysis):} High prealbumin suggests renal failure
  \item \textbf{Transthyretin genetic (TTR) testing:} If hereditary amyloidosis suspected (clinical: cardiomyopathy, neuropathy)
  \item \textbf{Pyrophosphate/C-PYP scintigraphy (PYP scan):} For transthyretin amyloid cardiomyopathy diagnosis
\end{itemize}

\section*{References}
\bibliographystyle{unsrtnat}
\bibliography{references}

\clearpage
\section*{Regulatory Status and Authorization Date}
This chapter describes a test category, not a specific commercial product. FDA, CDSCO, EU IVDR, and MHRA approval or registration dates therefore cannot be assigned without the assay manufacturer, product name, instrument, and intended use. For laboratory-developed tests, an individual agency approval date may not exist. See the Preface for the interpretation of regulatory status.

\clearpage
